% !TEX root = ../main.tex
\section{Modelo de dominio}\label{sec:dominio}

El modelo de dominio representa los conceptos del juego tal como los entiende un jugador, con sus atributos y asociaciones, sin operaciones ni decisiones de tecnología \citar{larman2003uml}. Los conceptos se obtuvieron de las reglas de negocio (sección~\ref{sec:reglas}) y de los casos de uso (sección~\ref{sec:casos}). La consola, el reloj del sistema y el almacenamiento en disco no forman parte del modelo, porque no son conceptos del juego sino medios técnicos.

\diagrama{modelo_dominio}{Modelo de dominio del juego de dominó}{fig:modelo-dominio}

\subsection{Conceptos del dominio}

La Tabla~\ref{tab:conceptos} describe cada concepto y la regla de negocio que lo origina. Los atributos precedidos de una barra en el diagrama son derivados: se calculan a partir de otros datos y no se registran por separado.

{\tablacompacta
\begin{longtable}{@{}L{3.2cm}L{8.4cm}L{2.9cm}@{}}
  \caption{Conceptos del modelo de dominio}\label{tab:conceptos}\\
  \toprule
  \cab{Concepto} & \cab{Descripción} & \cab{Origen} \\ \midrule \endfirsthead
  \toprule
  \cab{Concepto} & \cab{Descripción} & \cab{Origen} \\ \midrule \endhead
  \bottomrule \endlastfoot
  \texttt{Partida} & Encuentro configurado entre un jugador humano y de una a tres CPU. Registra el tipo de dominó, el límite de tiempo por turno y el estado. & CU-01, RN-01, RN-02, RN-04 \\ \filasep
  \texttt{Jugador} & Participante de la partida, identificado por su alias y su posición en la mesa. Se especializa en \texttt{JugadorHumano} y \texttt{JugadorCpu}. & RN-02 \\ \filasep
  \texttt{JugadorCpu} & Jugador automático con un nivel de dificultad: Principiante, Intermedio o Avanzado. & RN-02 \\ \filasep
  \texttt{Mano} & Fichas que posee un jugador y que los demás no ven mientras la partida está en curso. & RN-03, RN-09 \\ \filasep
  \texttt{Ficha} & Pieza con dos valores. Sus puntos y su condición de doble se derivan de esos valores. & RN-01, RN-05, RN-09 \\ \filasep
  \texttt{Mesa} & Cadena de fichas colocadas, con dos extremos abiertos cuyos valores se derivan de las fichas de los extremos. & RN-06 \\ \filasep
  \texttt{FichaColocada} & Ficha puesta en la mesa, con el extremo en el que se colocó y su orden. & RN-06 \\ \filasep
  \texttt{FichasParaRobar} & Fichas que quedan después del reparto y de las que roba el jugador sin ficha para lanzar. & RN-03, RN-07 \\ \filasep
  \texttt{Jugada} & Registro de cada acción del juego: colocar ficha, robar, pasar o perder el turno por tiempo. Conserva el orden en que ocurrieron. & RN-04, RN-10, RN-11 \\ \filasep
  \texttt{Resultado} & Forma en que terminó la partida (dominó o juego cerrado) y el ganador, si lo hay. & RN-08, RN-09 \\ \filasep
  \texttt{PuntajeFinal} & Puntos y cantidad de fichas que revela cada jugador cuando el juego se cierra. & RN-09 \\
\end{longtable}}

\subsection{Valores enumerados}

Los atributos con un conjunto cerrado de valores se representan como enumeraciones, de modo que las restricciones de la configuración y del estado quedan expresadas en el propio modelo.

\diagrama[0.85\textwidth]{enumeraciones}{Enumeraciones del modelo de dominio}{fig:enumeraciones}

\begin{table}[H]
\caption{Enumeraciones del dominio}\label{tab:enumeraciones}
\tablacompacta
\begin{tabularx}{\textwidth}{@{}L{3.4cm}Y@{}}
  \toprule
  \cab{Enumeración} & \cab{Valores} \\ \midrule
  \texttt{TipoDomino} & Doble6, Doble9 \\ \filasep
  \texttt{NivelDificultad} & Principiante, Intermedio, Avanzado \\ \filasep
  \texttt{LimiteTiempo} & SinLimite, Segundos10, Segundos20, Segundos30 \\ \filasep
  \texttt{EstadoPartida} & Preparada, Salida, EnCurso, Terminada \\ \filasep
  \texttt{Extremo} & Izquierdo, Derecho \\ \filasep
  \texttt{TipoJugada} & ColocarFicha, Robar, Pasar, TiempoAgotado \\ \filasep
  \texttt{FormaDeCierre} & Domino, JuegoCerrado \\
  \bottomrule
\end{tabularx}
\end{table}

\subsection{Restricciones del modelo}

\begin{itemize}
  \item Cada ficha se encuentra en un solo lugar a la vez: la mano de un jugador, la mesa o las fichas para robar.
  \item Una partida tiene exactamente un \texttt{JugadorHumano} y de uno a tres \texttt{JugadorCpu}.
  \item Cuando el juego se cierra, la partida pasa directamente a Terminada: se revelan las fichas de todos los jugadores y el resultado registra la forma de cierre \texttt{JuegoCerrado}. En la versión 1.0 el cierre se modelaba como un estado Cerrada previo, que se retiró porque no tenía reglas propias (sección~\ref{sec:bitacora}).
  \item Una partida terminada por dominó tiene ganador; una partida cerrada puede terminar empatada, en cuyo caso el resultado no tiene ganador.
\end{itemize}
